\section{Euler's integral and its effective calculus}
\label{sec:integrals}

\subsection{The improper integral is an algorithm}

Fix integers $A,R\ge1$ and a complex input with
\[
 \Rea z\ge1/A,
 \qquad |z|\le R.
\]
For positive $t$ use $t^{z-1}=\exp((z-1)\log t)$. Its modulus is
$t^{\Rea z-1}$, by the finite exponential identities. For
$0<\eta\le1$ and $T\ge1$,
\begin{align}
 \int_0^\eta |t^{z-1}e^{-t}|\dd t
   &\le A\eta^{1/A},\label{eq:gamma-low-tail}\\
 \int_T^\infty |t^{z-1}e^{-t}|\dd t
   &\le 2^{R+1}R!e^{-T/2}.\label{eq:gamma-high-tail}
\end{align}
The first inequality follows by integrating $t^{1/A-1}$. For the
second, $t^{\Rea z-1}\le t^R$ and the single term
$e^{t/2}\ge(t/2)^R/R!$ gives
$t^Re^{-t}\le2^RR!e^{-t/2}$. Initially these arguments apply between
finite positive endpoints; the displayed expressions are the certificates
for removing those endpoints.

Here are actual cutoffs for precision $p$. Put
\[
 D_R=2^{R+1}R!,\quad
 L=p+A+D_R+10,\quad
 \eta=2^{-AL},\quad T=2L.
\]
Then the sum of the two omitted tails is at most
$(A+D_R)2^{-L}<2^{-p-6}$. On $[\eta,T]$ a rational Lipschitz
majorant for the integrand is
\[
 C=(\eta^{-1}+T^R)\bigl((R+1)/\eta+1\bigr).
\]
Indeed its derivative is
$t^{z-1}e^{-t}((z-1)/t-1)$. The finite-calculus mesh of
\cref{sec:finite-calculus}, with $m\ge2^{p+6}(1+CT^2)$,
now specifies every evaluation. This gives a fixed primitive recursive
name $I(z)$ for the Euler integral.

\subsection{Identification with the finite-product gamma function}

For $N\in\N$, integration by parts on $[\eta,1]$ gives
\begin{equation}
 \int_0^1t^{z-1}(1-t)^N\dd t
 =\frac{N!}{z(z+1)\cdots(z+N)}.\label{eq:beta-polynomial}
\end{equation}
In detail, the $N=0$ antiderivative is $t^z/z$ and its lower boundary
has modulus at most $A\eta^{1/A}$. For $N>0$ integration by parts,
with antiderivative $t^z/z$, changes the integral to
$N/z$ times the same integral with $(z,N)$ replaced by $(z+1,N-1)$;
the boundary at $1$ is zero and the lower boundary tends to zero by
the same estimate. Iterating this finite relation proves
\eqref{eq:beta-polynomial}. To read the induction in $\ipra$, carry
the finite mesh and endpoint errors as parameters and decrease each
local error by the product bound \eqref{eq:product-budget}.

The substitution $t=s/N$ gives, for $N\ge1$,
\begin{equation}
 \Gamma_N(z)=\int_0^N s^{z-1}(1-s/N)^N\dd s.
 \label{eq:gamma-euler-integral}
\end{equation}
This formula directly identifies the two constructions with an explicit
rate. If $0\le s\le N/2$, the logarithm series proves
\[
 0\le-N\log(1-s/N)-s\le s^2/N,
\]
and hence
\[
 0\le e^{-s}-(1-s/N)^N\le e^{-s}s^2/N.
\]
On $[N/2,N]$ both kernels are ordered between $0$ and $e^{-s}$;
above $N$ only $e^{-s}$ remains. With
$C_R=1+2^{R+3}(R+2)!$, the elementary monomial bound used for
\eqref{eq:gamma-high-tail} therefore gives
\begin{equation}
 \boxed{\displaystyle
 |\Gamma_N(z)-I(z)|
 \le\frac{C_R}{N}+2^{R+1}R!e^{-N/4}\qquad(N\ge2).}
 \label{eq:integral-product-error}
\end{equation}
For the first term, split its integral at $1$: below $1$,
$s^{\Rea z+1}\le s$, and above $1$ it is at most $s^{R+2}$.
The resulting integral is bounded by $C_R$. The large tail begins at
$N/2$, proving the second term. For an entirely explicit choice set
$N=4(C_R+D_R+p+10)2^{p+4}$. Then each term in
\eqref{eq:integral-product-error} is below $2^{-p-3}$.

Both $I$ and $\Gamma$ are consequently equal to the same effective
Euler limit. We have proved
\begin{theorem}[Euler integral]
For a complex name with a supplied positive lower bound for its real part,
\[
 \boxed{\displaystyle\Gamma(z)=\int_0^\infty t^{z-1}e^{-t}\dd t.}
\]
The integral is absolutely convergent in the finite-tail sense of
\eqref{eq:gamma-low-tail}--\eqref{eq:gamma-high-tail}; it is not an
additional definition requiring an abstract completeness theorem.
\end{theorem}

\subsection{Positivity, modulus, and changes of variable}

For $x>0$, every finite integral summand is nonnegative and the interval
$1\le t\le2$ alone gives
\begin{equation}
 \Gamma(x)\ge1/18>0.\label{eq:gamma-positive}
\end{equation}
Here $t^{x-1}\ge1/2$ and $e^{-t}\ge e^{-2}>1/9$.
On $1/A\le x\le R$ we also have the useful upper bound
\begin{equation}
 \Gamma(x)\le A+2^{R+1}R!.\label{eq:gamma-real-bound}
\end{equation}
The triangle inequality on finite sums, followed by their tails, gives
$|\Gamma(\sigma+iy)|\le\Gamma(\sigma)$ for $\sigma>0$.
There is a stronger vertical-line product formula later in the report.

Integration by parts gives a second proof of
$\Gamma(z+1)=z\Gamma(z)$: the boundary $t^ze^{-t}$ has the same
explicit vanishing estimates. Also $\Gamma(1)=1$ follows from the
antiderivative $-e^{-t}$, so the factorial normalization can be proved
entirely from the integral as well.

Substitution on the truncated interval, followed by the already proved
endpoint bounds, gives
\begin{align}
 \Gamma(z)&=\int_{-\infty}^{\infty}e^{zu-e^u}\dd u,
 \label{eq:exponential-substitution}\\
 \Gamma(z)&=\int_0^1(-\log u)^{z-1}\dd u,
 \label{eq:unit-substitution}\\
 \Gamma(z)&=2c^z\int_0^\infty u^{2z-1}e^{-cu^2}\dd u
 \quad(c>0).\label{eq:quadratic-substitution}
\end{align}
For example \eqref{eq:exponential-substitution} uses
$t=e^u$ and the finite endpoints $u=\log\eta,\log T$;
\eqref{eq:unit-substitution} uses $t=-\log u$ and endpoints
$e^{-T},e^{-\eta}$; \eqref{eq:quadratic-substitution} uses
$t=cu^2$. Nonrational endpoints are evaluated with the local Lipschitz
budget. Thus the substitutions introduce no new convergence principle.
The more general positive-real scaling law is
\begin{equation}
 \int_0^\infty t^{z-1}e^{-ct^r}\dd t
 =\frac{\Gamma(z/r)}{r c^{z/r}}\quad(c,r>0,\ \Rea z>0),
 \label{eq:general-scaling}
\end{equation}
with rational lower and upper bounds for the positive parameters.
It follows from $u=ct^r$, whose derivative is explicit on every
truncated interval. Complex values of $c$ are not included in this
statement; they would require a specified complex power convention and
an additional argument.

\subsection{Differentiation under the integral, with tails}

For every natural number $j$, define the prospective derivative by
\[
 I_j(z)=\int_0^\infty
         t^{z-1}e^{-t}(\log t)^j\dd t.
\]
The additional logarithmic factors have explicit bounds. For
$0<t\le1$,
\[
 |\log t|^j\le j!(2A)^j t^{-1/(2A)},
\]
because $e^{v/(2A)}\ge(v/(2A))^j/j!$ for $v=-\log t$.
For $t\ge1$ use $(\log t)^j\le t^j$. Consequently
\begin{align}
 \int_0^\eta|t^{z-1}e^{-t}(\log t)^j|\dd t
 &\le j!(2A)^{j+1}\eta^{1/(2A)},\label{eq:derivative-low-tail}\\
 \int_T^\infty|t^{z-1}e^{-t}(\log t)^j|\dd t
 &\le2^{R+j+1}(R+j)!e^{-T/2}.\label{eq:derivative-high-tail}
\end{align}
These formulas give primitive recursive cutoff lengths uniformly in $j$.
On the remaining compact interval expand $e^{h\log t}$ to any fixed
degree, using \eqref{eq:exp-tail}. For example the first-order
remainder is at most
\[
 \tfrac12 |h|^2|\log t|^2e^{|h\log t|}.
\]
For $|h|\le1/(4A)$ this is integrable with exactly the same kind of
explicit bounds, using the smaller half-plane margin $1/(2A)$ and
radius $R+1$. Integrating the finite Taylor remainder proves
$I_j'=I_{j+1}$. Thus
\begin{theorem}[All integral derivatives]
For every $j\in\N$ and every certified input in the right half-plane,
\[
 \Gamma^{(j)}(z)
   =\int_0^\infty t^{z-1}e^{-t}(\log t)^j\dd t.
\]
The finite degree and integration errors are controlled by
\eqref{eq:derivative-low-tail}--\eqref{eq:derivative-high-tail} and
the compact Taylor remainder. In particular the statement is not an
unqualified exchange of an infinite limit with differentiation.
\end{theorem}

\subsection{Incomplete gamma as an additional finite-integral consequence}

For $x>0$ and $\Rea z>0$, put
\[
 \gamma_{\ell}(z,x)=\int_0^x t^{z-1}e^{-t}\dd t,
 \qquad \Gamma_u(z,x)=\int_x^\infty t^{z-1}e^{-t}\dd t.
\]
The subscript distinguishes the lower function from Euler's constant.
Additivity of finite sums and integration by parts prove
\begin{align*}
 \gamma_{\ell}(z,x)+\Gamma_u(z,x)&=\Gamma(z),\\
 \gamma_{\ell}(z+1,x)&=z\gamma_{\ell}(z,x)-x^ze^{-x},\\
 \Gamma_u(z+1,x)&=z\Gamma_u(z,x)+x^ze^{-x}.
\end{align*}
For a natural number $n$, induction from $\Gamma_u(1,x)=e^{-x}$ gives
\[
 \boxed{\displaystyle
 \Gamma_u(n+1,x)=n!e^{-x}\sum_{j=0}^n\frac{x^j}{j!}.}
\]
The lower function is $n!$ minus this expression. These formulas extend
to $x=0$ by the explicit small-endpoint estimate, with values $0$ and
$\Gamma(z)$ respectively; no decision whether a general name is zero
is needed in the continuity statement.
